\chapter{Conclusion}
\label{chapter:conclusion}
\thispagestyle{myheadings}

This thesis provides new perspectives on stochastic optimization through online-to-batch style techniques, with a focus on settings that are common in practical machine learning training such as differential privacy, non-convexity, and non-smoothness. Across Chapters 2--4, the main message is that carefully designed reductions can transfer algorithmic structure from online learning to challenging stochastic optimization problems while preserving desirable complexity guarantees.

Chapter 2 establishes a differentially private online-to-batch conversion for stochastic optimization with convex and smooth losses. The conversion framework transforms any online convex optimization method with $O(\sqrt{T})$ regret into an $\epsilon$-differentially private algorithm achieving the optimal rate $\tilde O(1/\sqrt{T} + \sqrt{d}/(\epsilon T))$ in linear time. Beyond this baseline guarantee, the framework supports adaptive variants whose rates depend on unknown variance or parameter norms, showing that private optimization can systematically inherit adaptivity from the online learning literature.

Chapter 3 extends this line to private non-convex and non-smooth optimization. The proposed zeroth-order conversion guarantees $(\alpha,\alpha\rho^2/2)$-R\'enyi differential privacy and returns a $(\delta,\epsilon)$-stationary point under near-optimal sample complexity. A key takeaway is that privacy can be obtained with essentially no extra asymptotic cost in certain regimes, even when smoothness is unavailable. This closes an important gap in private optimization by providing the first theoretical guarantees for the non-smooth non-convex setting.

Chapter 4 moves from differential privacy to non-smooth non-convex stochastic optimization and introduces Exponentiated Online-to-Non-Convex (Exp-O2NC). Using exponentially distributed random scaling and exponentiated regularization, Exp-O2NC achieves the optimal complexity guarantees in this setting. Pairing the reduction with an unconstrained online mirror descent subroutine recovers SGDM (with additional random scaling) and provides the first optimal convergence guarantees for this popular optimizer in the non-smooth non-convex setting. More broadly, this chapter shows that practical optimizers can be thought of as instances of the online-to-non-convex reduction, bridging training heuristics to online learning theory.

Taken together, the online-to-batch style reduction techniques in these chapters help analyze and establish theoretical guarantees in difficult yet practically common settings, including differential privacy, non-convexity, and non-smoothness. They can also be effective tools for designing optimizers: combining them with online learning algorithms allows us carry over useful properties from the rich literature of online learning, such as optimistic updates and adaptivity.

Several directions remain open. Regarding the DP SCO problem, our algorithm relies on the smoothness assumption to achieve theoretical guarantees. Reducing or even removing smoothness-related restrictions would make our algorithm more generally applicable. Regarding the private non-convex problem, open directions include designing more efficient estimators with lower oracle complexity and understanding the matching lower bounds in this setting. On the exponentiated O2NC side, a promising avenue is adaptive variants that connect more directly to Adam-style methods while retaining non-asymptotic guarantees. Progress on these questions would further strengthen the unifying role of online-learning-based reductions in optimization theory and practice.
